\documentclass[12pt]{article}
\usepackage[margin=1in]{geometry}
\usepackage{enumitem}
\usepackage{titlesec}
\usepackage{parskip}
\usepackage{graphicx}
\usepackage{xcolor}
\usepackage{hyperref}
\usepackage{fontenc}

\title{\textbf{Harmel and Janda 1994 RG Notes}\\
\large Harmel, Robert, and Kenneth Janda. ``An Integrated Theory of Party Goals and Party Change'' [1994].\\
\textit{Journal of Theoretical Politics} 6(3): 259-287.}
\author{}
\date{}

\begin{document}

\maketitle

Party change results from one of three sources:

1) leadership change

2) a change of dominant faction within the party

3) an external stimulus for change

Very similar to the EVL framework from Clark et al 2017, at least in highlighting \textit{why} the scenario at hand might change

Previous literature often suggests volatility or exogenous decline; here, suggested that a series of behaviors, changes, and actions are necessary to meaningfully influence party direction

Three primary premises of the theory:

1) Each party has a primary goal which can vary both across party and across time. (This is a quick link to Policy/Office/Votes by Strom 1990 and Muller/Strom 1999)

2) Dramatic change occurs after shocks: parties are conservative organizations in that they will not substantially change without good cause

3) External shocks = ``external stimuli that impact on the party's primary goal'' -- this can change across parties based on their goal. For example, regime change may affect policy-driven parties much more than ideology-driven parties or electorally-driven parties

\textbf{Adds a fourth goal beyond Strom 1990: ``Implementing Party Democracy'', or ``democracy-seeking.'' Parties can be motivated by a desire to be democratic themselves, though this is not necessarily a given.} (This is a huge contribution to the reading list even though it is only briefly mentioned here)

\end{document}
